\begin{pastexam}[label={px:nccu-114-dsa-08}]{政大資科所 114 學年度 第 8 題}
（單選；每題 4 分）

(4\%) If an array is used to implement a max-heap, where index 1 (the first position) of the array
stores the root, what is the maximal index $i$ that stores the 3rd ranked element?

A.~3 \qquad B.~4 \qquad C.~5 \qquad D.~6 \qquad E.~7
\end{pastexam}

\begin{pastsoln}
答案是 E。 「3rd ranked element」指第 $3$ 大的元素，假設元素互不相同、堆積至少有 $7$ 個元素。

最大堆積（資結書第~8.1~節）裡，每個節點都比它的子孫大。 所以一個節點的每一個祖先
都比它大。 由此：
\begin{itemize}
\item 第 $3$ 大的元素最多只有 $2$ 個元素比它大，所以它最多有 $2$ 個祖先，深度 $\le 2$（根的深度
是 $0$）。 深度 $\le 2$ 的節點索引是 $1$ 到 $7$，所以索引 $\le 7$。
\item 索引 $7$ 做得到：讓第 $2$ 大的放在索引 $3$（根的右子節點），第 $3$ 大的放在索引 $7$（索引
$3$ 的右子節點）。 例如
\[
A[1..7] = \langle 10, 5, 9, 1, 2, 3, 8 \rangle ,
\]
每個父節點都比子節點大，是合法的最大堆積，第 $3$ 大的 $8$ 在索引 $7$。
\end{itemize}
所以最大索引是 $7$。 （$n < 7$ 時答案是 $n$，例如 $n = 4$ 時第 $3$ 大可以在索引 $4$；選項是以
堆積夠大為前提。）

逐項檢查：A、B、C、D 的索引都放得下第 $3$ 大的元素，但都不是最大的；E 的 $7$ 是深度 $2$ 的最後
一個位置，上面已經證明不可能更大，也舉出了做得到的例子。 答案 E。
\end{pastsoln}
